\section{Tổng quan sản phẩm (Product Overview)}
\label{sec:srs_product_overview}

\subsection{Mô tả hệ thống (System Description)}
\label{subsec:srs_system_description}

\subsection{Sơ đồ ngữ cảnh (Context Diagram)}
\label{subsec:srs_context_diagram}

\subsection{Quy trình nghiệp vụ tổng quát (Business Process)}
\label{subsec:srs_business_process}

\section{Yêu cầu người dùng (User Requirements)}
\label{sec:srs_user_requirements}

\subsection{Danh sách Actor (Actor List)}
\label{subsec:srs_actor_list}

Dưới đây là mô tả chi tiết vai trò của các đối tượng sử dụng (Actor) tham gia vào hệ thống Manga Creation Workflow and Publishing Management System:

\begin{itemize}
    \item \textbf{Admin (Quản trị viên):} Là người quản trị hệ thống, chịu trách nhiệm quản lý tài khoản người dùng, phân quyền truy cập và đảm bảo hệ thống hoạt động ổn định. Admin có quyền kiểm soát các chức năng quản trị và theo dõi hoạt động chung của hệ thống.
    \item \textbf{Mangaka (Tác giả):} Là tác giả chính của tác phẩm Manga, chịu trách nhiệm xây dựng nội dung, quản lý quá trình sáng tác và phối hợp với các Assistant để hoàn thiện tác phẩm trước khi gửi kiểm duyệt cho Tantou Editor.
    \item \textbf{Assistant (Trợ lý):} Là người hỗ trợ Mangaka thực hiện các công việc được phân công (vẽ nền, tô âm ảnh, đi nét hoặc chỉnh sửa chi tiết đồ họa...) theo sự giao việc của Mangaka để hoàn thiện từng trang truyện.
    \item \textbf{Tantou Editor (Biên tập viên phụ trách):} Là biên tập viên trực tiếp theo dõi quá trình phát triển của các Series Manga được phân công. Tantou Editor chịu trách nhiệm kiểm duyệt nội dung, đưa ra nhận xét phản hồi và hỗ trợ Mangaka nâng cao chất lượng tác phẩm trước khi xuất bản.
    \item \textbf{Editorial Board (Hội đồng biên tập):} Là hội đồng biên tập chịu trách nhiệm đánh giá hiệu quả hoạt động của các Series Manga dựa trên dữ liệu bình chọn và kết quả phát hành. Hội đồng đưa ra các quyết định liên quan đến việc tiếp tục xuất bản, thay đổi lịch phát hành hoặc ngừng phát hành Series.
\end{itemize}

\subsection{Use Case Diagram}
\label{subsec:srs_use_case_diagram}

\subsection{Bảng tóm tắt Use Case (Use Case Summary)}
\label{subsec:srs_use_case_summary}

\subsection{Yêu cầu của từng nhóm người dùng}
\label{subsec:srs_user_specific_requirements}
